\chapter{Estado del arte}
\label{cap:estado}

\section{Firmas génicas en cáncer de pulmón}
\label{sec:firmas_estado}

La búsqueda de firmas génicas con valor diagnóstico o pronóstico en cáncer
de pulmón cuenta con una literatura extensa. Beer et
al.~\cite{beer2002gene} describieron en 2002 una de las primeras firmas
pronósticas para adenocarcinoma de pulmón mediante \emph{microarrays},
identificando un panel de aproximadamente 50 genes cuya expresión permitía
estratificar pacientes por riesgo de muerte a corto plazo. En los años
siguientes, distintos grupos han propuesto firmas para diagnóstico
histológico, pronóstico global, respuesta a tratamiento con inhibidores de
EGFR, presencia de mutaciones en KRAS o ALK, y probabilidad de recurrencia
tras cirugía.

Sin embargo, la revisión sistemática de Subramanian y Simon~\cite{subramanian2010gene}
reveló que la gran mayoría de firmas propuestas no superaba una validación
independiente rigurosa. Los sesgos habituales incluyen: (i) uso de la misma
cohorte para descubrimiento y validación mediante partición aleatoria en
lugar de LODO; (ii) selección de genes en el conjunto completo de datos
seguida de validación cruzada, incurriendo en la falacia de las
particiones (\emph{data leakage}); (iii) omisión de análisis de composición
tisular, con firmas que en realidad reflejan la proporción variable de
estroma o células inmunes; y (iv) informar únicamente el AUC sin
considerar la capacidad de calibración del umbral de decisión entre
cohortes.

Estos problemas metodológicos explican, al menos en parte, por qué el
número de firmas génicas trasladadas a la práctica clínica en cáncer de
pulmón es reducido. Uno de los ensayos comerciales con mayor difusión es la
firma pronóstica \emph{Pervenio Lung RS}, basada en un panel de 14 genes
medidos por PCR cuantitativa~\cite{kratz2012practical}; otro es
\emph{Oncotype DX Lung} de Genomic Health, con un panel similar. Ambos
mantienen un uso limitado y no han desplazado a los marcadores
inmunohistoquímicos clásicos.

En el ámbito del diagnóstico histológico, la Organización Mundial de la
Salud (OMS) publica desde 2015 recomendaciones para el uso combinado de
marcadores inmunohistoquímicos: TTF-1/NKX2-1 y napsina A como positivos
para adenocarcinoma, y CK5/6, p40 (isoforma de TP63) y DSG3 como positivos
para carcinoma escamoso~\cite{travis2015who}. La correspondencia entre
estos marcadores IHC y su expresión transcripcional constituye una
referencia natural para validar la interpretabilidad de una firma génica.

\section{Integración multi-cohorte y meta-análisis}
\label{sec:metaanalisis}

La integración de múltiples cohortes de expresión ha dado lugar a un cuerpo
metodológico específico. Los enfoques principales se agrupan en dos
familias:

\begin{itemize}
  \item \textbf{Corrección de efecto lote:} normaliza las cohortes
  conjuntamente para eliminar la varianza atribuible al origen técnico.
  ComBat~\cite{johnson2007adjusting} y sus variantes bayesianas empíricas
  son las herramientas más difundidas. El riesgo de estos métodos es
  eliminar señal biológica genuina cuando el efecto lote está confundido
  con la variable de interés.
  \item \textbf{Meta-análisis por consenso:} normaliza cada cohorte por
  separado, calcula el estadístico de interés (por ejemplo, el tamaño de
  efecto \emph{d} de Cohen para la diferencia de medias entre grupos), y
  combina los resultados con criterios de consenso (mismo signo, misma
  magnitud mínima) en varias cohortes. Este enfoque es más conservador y
  permite distinguir señal reproducible de artefacto técnico.
\end{itemize}

En el contexto de este trabajo se adopta la segunda familia, con validación
externa Leave-One-Dataset-Out (LODO). En LODO, el modelo se entrena con
todas las cohortes menos una, y se evalúa en la retirada; el proceso se
itera dejando fuera cada cohorte por turno. Esta partición es más
conservadora que la validación cruzada estándar porque respeta la
estructura de agrupación de las muestras: dos muestras de la misma cohorte
comparten sesgo técnico, y evaluarlas juntas sobre-estima el rendimiento.

Bernau et al.~\cite{bernau2014crossvalidation} formalizaron el marco de
validación cross-study en clasificadores oncológicos. Trabajos posteriores
han aplicado LODO a firmas de cáncer de mama~\cite{waldron2014comparative}
y de pulmón~\cite{shukla2017cross}, mostrando reducciones consistentes del
AUC entre 5 y 15 puntos cuando se pasa de validación interna a externa.
Esto justifica adoptar LODO como estándar mínimo para reclamar
replicación.

\section{Aprendizaje automático en transcriptómica}
\label{sec:ml_estado}

Los datos transcriptómicos plantean un reto clásico para el aprendizaje
automático: alta dimensionalidad ($p \gg n$, con $p$ del orden de $10^4$
genes y $n$ del orden de $10^2$ muestras). Las estrategias habituales
combinan reducción de dimensionalidad, regularización y selección de
características.

\begin{itemize}
  \item \textbf{Regresión logística con regularización L1} (LASSO): la
  penalización $\ell_1$ fuerza a coeficientes exactos a cero, seleccionando
  implícitamente un subconjunto reducido de genes. Es la estrategia más
  usada cuando se busca interpretabilidad y un panel mínimo.
  \item \textbf{Random Forest:} ensamble de árboles de decisión con
  \emph{bagging} y selección aleatoria de variables en cada división. No
  fuerza esparsidad pero produce un ranking de importancia útil como
  complemento a LASSO.
  \item \textbf{Máquinas de vectores soporte} (SVM): eficaces en
  dimensiones altas con kernels lineales o RBF, aunque con menor
  interpretabilidad de los pesos.
  \item \textbf{Redes neuronales profundas:} en los últimos años se han
  aplicado a datos ómicos, pero requieren tamaños muestrales que suelen
  quedar fuera del alcance de estudios de cohortes públicas.
\end{itemize}

La elección de LASSO como método principal en este trabajo se justifica por
tres razones: (i) produce un panel interpretable, útil clínicamente;
(ii) es determinista dado un \texttt{random\_state} fijo, lo que garantiza
reproducibilidad; y (iii) el hiperparámetro de regularización (C) admite
una calibración simple mediante validación cruzada anidada, sin requerir
las decisiones arquitectónicas propias de una red neuronal.

\section{Modelos de lenguaje en bioinformática}
\label{sec:llm_estado}

El uso de LLMs en bioinformática ha crecido rápidamente desde la aparición
de GPT-4 en 2023. Las aplicaciones documentadas incluyen anotación
funcional de genes, extracción de información de literatura biomédica,
generación de código para análisis y curación de metadatos
clínicos~\cite{tinn2023bioclinicalbert, thirunavukarasu2023llm}. Un trabajo
particularmente relevante para esta tesis es GEOparser~\cite{wang2024geoparser},
que emplea GPT-4 para extraer metadatos estructurados de GEO Series y
demuestra tasas de éxito superiores al 90\,\% en un conjunto de test
manual.

Los modelos abiertos (Llama~2, Llama~3, Mistral) han cerrado buena parte
de la brecha con los modelos propietarios en tareas de comprensión y
clasificación, especialmente en tamaños de 70\,B parámetros o
superiores~\cite{touvron2023llama}. La infraestructura de inferencia
acelerada por hardware específico (Groq LPU, Cerebras) permite servir estos
modelos con latencias de decenas de milisegundos por consulta, haciendo
viable su integración en \emph{pipelines} de análisis a escala.

En el contexto de este trabajo, el LLM cumple una única función bien
delimitada: leer la descripción textual de una muestra de GEO y devolver una
etiqueta canónica del grupo experimental. Esta tarea es paradigmática de
clasificación de texto y no requiere generación abierta, lo que reduce los
riesgos de alucinación observados en tareas de razonamiento libre. El uso
del LLM se restringe además a la fase de curación de datos, sin
participación en el análisis estadístico ni en la selección de
biomarcadores, evitando así que el conocimiento paramétrico del modelo
sobre cáncer de pulmón introduzca sesgos en los resultados.

\section{Ausencias en la literatura}
\label{sec:huecos}

La revisión anterior identifica varios huecos que este TFM aborda de
manera integrada:

\begin{enumerate}
  \item \textbf{Ausencia de \emph{frameworks} extremo a extremo} que
  combinen curación automatizada con LLM, integración multi-cohorte,
  clasificación supervisada con validación externa y verificación cruzada
  contra marcadores clínicos, todo ello en un pipeline reproducible.
  \item \textbf{Falta de firmas con replicación estricta} en varias
  cohortes independientes. Los trabajos publicados suelen reportar
  resultados sobre una única cohorte de validación o sobre particiones
  aleatorias de la misma cohorte de descubrimiento.
  \item \textbf{Desconexión entre firmas transcriptómicas y marcadores
  clínicos.} La mayoría de firmas propuestas no verifican explícitamente
  hasta qué punto recuperan marcadores ya en uso en la clínica, ni
  explican qué genes de su panel corresponden a marcadores conocidos y
  cuáles son novedosos.
\end{enumerate}

El framework \framework\ se diseña como respuesta a estos tres huecos:
proporciona una tubería end-to-end reproducible, aplica LODO como criterio
mínimo de validación y contrasta explícitamente la firma resultante con el
panel diagnóstico IHC de la OMS.
